Capability scaling in modern large language models (LLMs) has been predominantly driven by parameter expansion, with Mixture-of-Experts (MoE) architectures emerging as the de facto standard for achieving this at a tractable cost~\cite{liu2024deepseek, yang2025qwen3, zeng2025glm}. The efficacy of MoE stems from a structural decoupling, in which activating only a sparse subset of experts per token disaggregates total parameter count from per-token active compute and enables parameter growth without commensurate FLOP inflation. However, scaling total parameters is not the sole avenue for enhancing capability per unit of training cost. An orthogonal paradigm investigates how to extract greater computational depth from a strictly bounded parameter budget. Among these approaches, block-recurrent or looped architectures offer a compelling solution by executing a shared block of layers over multiple iterations, thereby increasing depth without introducing new parameters~\cite{dehghani2018universal, lan2019albert, bae2025relaxed, mohtashami2025cotformer, bae2026mixture, geiping2026scaling}.

While block-recurrent architectures successfully increase computational depth, existing dense loops~\cite{zhu2025scaling, zeitoun2026hyperloop, jeddi2026loopformer} suffer from a fundamental structural flaw: parameter count and per-token FLOPs remain tightly coupled. Iterating a dense block multiplies compute without altering parameters, which makes it impossible to simultaneously match a non-looped baseline on both axes. MoE architectures~\cite{cai2025survey}, which naturally decouple total parameters from active compute, offer a theoretical solution. However, naively integrating loops into an MoE backbone introduces two critical new challenges. Representationally, weight sharing across iterations imposes structural symmetry that severely restricts the model's ability to evolve token states progressively. Structurally, it induces asymmetrical active parameter expansion: attention parameters are reused uniformly across iterations, whereas tokens dynamically route to different experts at each step. As a consequence, the active attention-to-FFN parameter ratio $\rho$ deviates substantially from the operating point $\rho^\star$ of well-tuned baselines.

To fully unlock the potential of iterative sparse computation, we introduce \textbf{\ourmodelnospace}, a novel block-recurrent MoE language model designed to overcome both limitations. \ourmodel adopts a streamlined sandwich layout in which the core loop body evolves through pure recursion. To break the structural symmetry inherent in weight sharing, we introduce \textbf{IterAdaLN}, a token-level modulation scheme inspired by adaptive normalization in conditional generation~\cite{perez2018film, peebles2023scalable}. Operating strictly at a token-level granularity, IterAdaLN dynamically generates affine parameters jointly from the iteration index and the per-token hidden state. Replacing static RMSNorm~\cite{zhang2019root} with this token-level conditioning ensures sufficient differentiation across iterations. It also allows the representational trajectory to remain highly expressive without requiring iteration-specific prefix re-injection.

Crucially, \ourmodel leverages the inherent flexibility of the MoE framework to resolve the asymmetrical expansion problem. To maintain the optimal ratio $\rho^\star$, we adopt a principled balancing strategy that expands the Q/KV LoRA ranks of Multi-head Latent Attention (MLA)~\cite{liu2024deepseek} and correspondingly reduce the expert hidden dimensions. This decoupling allows independent adjustment of attention and FFN capacities. We then recover the baseline's total parameter count $N^\star$ by scaling the routed expert count. Together, these designs turn the MoE structure into the exact degree of freedom needed to construct a rigorously controlled looped architecture.

% Combining token-modulated block recurrence with a principled balancing strategy, \ourmodel provides the first rigorously matched evaluation of looped architectures. Under strictly matched total parameters, per-token compute, and active parameter ratios, and even with slightly fewer active parameters per token, \ourmodel improves over a strong non-loop MoE baseline on $8$ out of $9$ benchmarks by up to roughly 2.3 points and lifts the overall average by more than 1 point. The improvements are most pronounced on tasks demanding reasoning ability. Notably, on reasoning and mathematics benchmarks \ourmodel remains competitive with OLMoE-1B-7B~\cite{muennighoff2025olmoe}, a model with over twice the total parameter count. Because our evaluation strictly isolates the architectural design from parameter or FLOP inflation, it provides definitive evidence that the performance gains stem from the intrinsic superiority of the block-recurrent architecture. Internal analysis further corroborates this finding, showing that the loop body self-organizes into a back-loaded update schedule, where later iterations drive the majority of per-token differentiation. Together, these results validate \ourmodel as a highly effective paradigm for scaling reasoning capabilities within a fixed capacity budget.

Our main contributions are summarized below:
\begin{itemize}[leftmargin=*]
\item \textbf{Novel Architecture}: We propose \textbf{\ourmodelnospace}, which couples sparse expert routing with iterative weight-shared computation through two key designs: \textbf{IterAdaLN}, a token-level iteration-conditioned modulation that breaks weight-sharing symmetry, and a capacity-balancing strategy that maintains the attention-to-FFN active-parameter ratio of well-tuned non-looped MoEs.
\item \textbf{Strictly Controlled Evaluation}: We provide the first head-to-head comparison of a looped architecture against a Vanilla MoE under rigorously matched total parameters, per-token FLOPs, and active sublayer ratios, cleanly isolating the architectural contribution from parameter or compute inflation.
% \item \textbf{Consistent Empirical Gains}: Under this matched setting, LoopMoE outperforms the Vanilla MoE on 8 of 9 benchmarks while using fewer physical active parameters per token. The average improvement exceeds 1.0 point, with the largest single-task gain reaching 2.3 points. On reasoning and mathematics, LoopMoE further remains competitive with OLMoE-1B-7B~\cite{muennighoff2025olmoe}, despite the latter having over twice the total parameter count.
\item \textbf{Consistent and Scale-Robust Empirical Gains}: Under this matched setting at 3B scale, LoopMoE outperforms the Vanilla MoE on 8 of 9 benchmarks while accessing fewer physical active parameters per token, with the average improvement exceeding 1.0 point. On reasoning and mathematics, LoopMoE further remains competitive with OLMoE-1B-7B, despite the latter having over twice the total parameter count. At the 9B scale, LoopMoE continues to outperform the matched Vanilla MoE at an early-training checkpoint, indicating that the architectural benefit is preserved at a larger scale rather than being a small-scale artifact.
\end{itemize}